\chapter{The angle and slope of a graphical flow}\label{sec:graphical-flow-angle}
Ecker--Huisken's gradient estimates begin with the reciprocal vertical angle; see \cite{EH89}, Lemma 3.1 and Corollary 3.2, printed p.~458. We derive the local identities first and apply the maximum principle on a compact periodic graph. This avoids any implicit maximum-principle assumption at spatial infinity. The flat formulas below also hold for curves, so here $n\geq1$.

\begin{definition}[Periodic graphs and normal material time]\label{def:periodic-graph-flow}
Let $\Lambda\subset\mathbb R^n$ be a full-rank lattice, $\mathbb T^n=\mathbb R^n/\Lambda$ its flat torus, and let $u:\mathbb T^n\times I\to\mathbb R$ be smooth. A periodic graphical mean curvature flow is the graph of $u$ in the flat product $\mathbb T^n\times\mathbb R$, with $u$ satisfying Proposition \ref{prop:graphical-mcf}. Equivalently, its lift is a periodic entire graph in Euclidean space. With $e$ the upward parallel unit vector, set
\[
 \nu=\frac{(-Du,1)}{v},\qquad v=\sqrt{1+|Du|^2},\qquad w=\langle\nu,e\rangle=v^{-1}>0.
\]
We use $A=D\nu$ and $\vec H=-H\nu$. The symbol $\partial_t^\perp$ denotes differentiation in a parametrization with velocity $-H\nu$, and $\mathcal P=\partial_t^\perp-\Delta_{M_t}$. It is not automatically the derivative at a fixed graph coordinate. This is the angle convention of \cite{EH89}, Section 3.
\end{definition}

\begin{proposition}[Angle, reciprocal angle and their powers]\label{prop:graphical-angle-evolution}
On any oriented smooth flat-ambient flow with $w=\langle\nu,e\rangle>0$ and $v=w^{-1}$,
\[
 \mathcal Pw=|A|^2w,\qquad
 \mathcal Pv=-|A|^2v-2v^{-1}|\nabla v|^2.
\]
For every real $p$,
\[
 \mathcal P(v^p)=-p|A|^2v^p-p(p+1)v^{p-2}|\nabla v|^2,
 \qquad
 \mathcal P\log v=-|A|^2-v^{-2}|\nabla v|^2.
\]
These are pointwise identities wherever the angle is positive, independently of compactness. The reciprocal-angle equation is \cite{EH89}, Lemma 3.1.
\end{proposition}
\begin{proof}
Proposition \ref{prop:gauss-map-tension} already proves $\Delta_M\nu=\nabla H-|A|^2\nu$. Proposition \ref{prop:normal-variation-formulas} gives $\partial_t^\perp\nu=\nabla H$. Subtract and pair with the parallel vector $e$ to obtain the equation for $w$. These local identities also hold on a flat quotient and in dimension one. The scalar chain rule for $\mathcal P$ is
$\mathcal P\Phi(w)=\Phi'(w)\mathcal Pw-\Phi''(w)|\nabla w|^2$.
Take $\Phi(w)=w^{-1}$ and use $|\nabla w|^2=v^{-4}|\nabla v|^2$. This proves the equation for $v$. Apply the same rule to $v^p$; the two gradient coefficients combine to $-2p-p(p-1)=-p(p+1)$. Taking $\Phi(v)=\log v$ gives the last formula directly. No sign restriction on $p$ is needed because $v$ is positive.
\end{proof}

\begin{lemma}[The drift introduced by graph coordinates]\label{lem:graphical-material-drift}
For the parametrization $F(x,t)=(x,u(x,t))$, its velocity decomposes as
\[
 F_t=-H\nu+T,\qquad T=u_t e^T=-Hv e^T,
\]
where $e^T=e-w\nu$. If $f$ is a geometric scalar on the moving graph, then
\[
 \partial_t^{\rm graph}f=\partial_t^\perp f+\langle T,\nabla f\rangle.
\]
Thus an equation involving $\mathcal P+b\cdot\nabla$ becomes, in graph coordinates, the equation with operator
$\partial_t^{\rm graph}-\Delta+(b-T)\cdot\nabla$.
\end{lemma}
\begin{proof}
The graph equation says $u_t=-Hv$. Since $e=w\nu+e^T$, its velocity $u_te$ has normal part $u_tw\nu=-H\nu$ and the asserted tangential part. Removing that tangential velocity by a time-dependent parameter change gives a normal parametrization locally in time. The ordinary chain rule for a scalar under this parameter change gives the displayed derivative identity. On any compact smooth time slab of a periodic graph, $T$ and all coefficients are smooth and bounded. Hence the tangential term is a legitimate drift in the compact maximum principle, and vanishes at a spatial extremum of $f$. This explains why an angle calculation in normal time can be used for a graph-coordinate estimate.
\end{proof}

\begin{proposition}[Slope cannot grow on a periodic graph]\label{prop:periodic-graph-slope-bound}
For a smooth periodic graphical flow on $[s,T)$,
\[
 1\leq v(x,t)\leq V_s:=\max_{M_s}v\quad(s\leq t<T).
\]
Consequently $|Du(x,t)|\leq\max|Du(\cdot,s)|$, and the graphical equation has spatial ellipticity between $V_s^{-2}$ and $1$ throughout this interval. This is the periodic version of \cite{EH89}, Corollary 3.2; it does not assert global existence.
\end{proposition}
\begin{proof}
The reciprocal-angle equation has $\mathcal Pv\leq0$. By Lemma \ref{lem:graphical-material-drift}, its graph-coordinate form has only an additional smooth tangential drift. On any compact smooth slab $[s,t_1]\subset[s,T)$, the maximum principle therefore gives $\max v(t)\leq\max v(s)$. One can obtain the weak principle directly by applying the first interior maximum test to $v-V_s-\eta(t-s)$ and then letting $\eta\downarrow0$. Compactness supplies a maximum; there is no lateral boundary. The lower bound follows from the graph formula. The gradient bound and ellipticity are immediate from $v^2=1+|Du|^2$ and Proposition \ref{prop:graphical-mcf}.
\end{proof}

\begin{lemma}[The gradient of the reciprocal angle]\label{lem:graphical-slope-gradient}
Pointwise on a graphical hypersurface,
\[
 \nabla w=A(e^T),\qquad \nabla v=-v^2 A(e^T),\qquad
 |\nabla v|\leq v\sqrt{v^2-1}\,|A|.
\]
In particular $v^{-1}|\nabla v|\leq v|A|$. These estimates control the drift in the weighted curvature equation when slope and curvature are bounded, as in \cite{EH89}, Corollary 4.2.
\end{lemma}
\begin{proof}
For a tangent vector $X$, differentiate $w=\langle\nu,e\rangle$ to obtain $Xw=\langle AX,e^T\rangle$. Self-adjointness of $A$ identifies its gradient. Differentiate $v=w^{-1}$ for the second formula. Finally $|e^T|^2=1-w^2=1-v^{-2}$ and the operator norm of $A$ is at most its tensor norm. Hence $|\nabla v|\leq v^2|A|\sqrt{1-v^{-2}}$, which is the stated bound, including at points where $v=1$ or $A=0$.
\end{proof}
